% Chapter2/chapter2.tex -- Literature Review
% -----------------------------------------------------------------------------
% The literature review has two jobs:
%   1. Demonstrate that you understand the research landscape
%   2. Justify your design decisions by showing what others have tried
%
% Structure tip: don't just summarise papers one-by-one. Group them by theme
% (quantisation, attention optimisation, FPGA architectures) and synthesise
% them -- explain how they relate to each other and to YOUR work.

\chapter{Literature Review}
\label{ch:litreview}

\section{Transformer Model Architecture}
\label{sec:transformer_arch}

The transformer architecture, introduced by Vaswani et al.\ \cite{PLACEHOLDER},
forms the computational backbone of all modern LLMs including TinyLlama.
Unlike recurrent networks, transformers process sequences in parallel through
a mechanism called self-attention, which allows every token to directly attend
to every other token in the sequence. While this parallelism enables efficient
training on GPUs, it creates a challenging memory-access pattern for hardware
accelerators: the attention computation scales quadratically with sequence
length, and the weight matrices must be repeatedly loaded from external DRAM
during inference.

The TinyLlama-1.1B model \cite{PLACEHOLDER} uses the LLaMA-2 architecture with
grouped-query attention (GQA) and the SwiGLU activation function in place of
the standard ReLU. GQA reduces the number of key and value heads relative to
query heads, substantially decreasing the KV cache size during autoregressive
generation -- a particularly important optimisation for memory-constrained
hardware like the ZCU102.

\lipsum[2] % Replace with your actual survey of transformer variants

\section{FPGA Acceleration of Neural Networks}
\label{sec:fpga_acceleration}

FPGA-based neural network acceleration has been an active research area since
the mid-2010s. Early work such as \cite{PLACEHOLDER} demonstrated that FPGAs
could outperform CPUs on convolutional neural network (CNN) inference by
exploiting data-level parallelism and on-chip memory buffering. The key insight
is that FPGAs allow the designer to co-optimise the compute pipeline and the
memory hierarchy simultaneously -- something that is fixed in GPU silicon.

For transformer models specifically, the dominant challenge is memory bandwidth
rather than raw arithmetic throughput. A single forward pass of TinyLlama-1.1B
moves approximately 2.2\,GB of weight data from DRAM, yet performs only around
2.2\,GFLOP of arithmetic -- giving a compute-to-memory ratio of roughly
1\,FLOP/byte, which is far below what modern FPGAs can sustain at their
arithmetic peak. This memory-bound characteristic means that the primary
optimisation target is minimising redundant DRAM accesses through on-chip
weight buffering and tiling, rather than maximising multiply-accumulate
throughput.

\lipsum[3] % Replace with your survey of FPGA transformer acceleration papers

\section{Quantisation for Efficient Inference}
\label{sec:quantisation}

Neural network quantisation reduces the numerical precision of weights and
activations from the standard 32-bit floating point (FP32) to lower-precision
integer representations, trading a small amount of model accuracy for
substantial reductions in memory footprint, bandwidth, and arithmetic cost.
For FPGAs in particular, fixed-point integer arithmetic is significantly more
resource-efficient than floating-point: a single-precision floating-point
multiplier requires roughly 3$\times$ as many LUTs as an equivalent INT8
multiplier on Xilinx UltraScale+ devices \cite{PLACEHOLDER}.

Post-training quantisation (PTQ) at INT8 precision has become the dominant
approach for production LLM deployment, with works such as LLM.int8()
\cite{PLACEHOLDER} and SmoothQuant \cite{PLACEHOLDER} demonstrating that
8-bit quantisation of both weights and activations incurs less than 1\%
accuracy degradation on standard benchmarks when applied with appropriate
per-channel scaling factors.

\lipsum[4] % Replace with your survey of quantisation methods

\section{HLS-Based Hardware Design}
\label{sec:hls_design}

High-level synthesis tools bridge the gap between algorithmic descriptions in
C/C++ and synthesisable RTL, dramatically reducing the design iteration cycle
compared to manual HDL development. Vitis HLS, the primary tool used in this
project, translates annotated C++ functions into pipelined hardware modules
whose behaviour is controlled through pragma directives embedded in the source
code \cite{PLACEHOLDER}.

The most critical HLS directives for throughput optimisation are PIPELINE
(which enables concurrent execution of successive loop iterations, characterised
by the initiation interval II), UNROLL (which replicates loop body logic to
increase parallelism), and ARRAY\_PARTITION (which splits arrays into multiple
memory banks to eliminate memory-port bottlenecks). Prior work on HLS-based
transformer acceleration \cite{PLACEHOLDER} has shown that achieving II\,=\,1
on the innermost matrix multiplication loop is the single most impactful
optimisation, as it determines the fundamental throughput ceiling of the
entire pipeline.

\lipsum[5] % Replace with your survey of HLS-based ML accelerator papers

\section{Summary and Research Gap}
\label{sec:research_gap}

Existing FPGA accelerators for transformer inference have largely focused on
encoder-only models (BERT and its variants) for classification tasks, where
the sequence length is fixed and the computational graph is static. Accelerators
for autoregressive decoder-only models like TinyLlama face an additional
challenge: each generation step produces a single token and must update the KV
cache, creating a highly sequential data dependency that is difficult to
pipeline across time steps. To the best of the author's knowledge, no published
work has demonstrated a full end-to-end implementation of a 1B+ parameter
decoder-only LLM on the ZCU102 platform using Vitis HLS with INT8 quantisation.
This project addresses that gap.
